\centerline{\bf \Huge Мюрхед}

\begin{flushright}
        \scriptsize{листочек для фарма ауры}
\end{flushright}

Пусть даны переменные $x_1, x_2, \ldots, x_n$. \textit{Однородным симметрическим многочленом}, порождённым упорядоченным набором целых неотрицательных чисел $a_1 \geqslant a_2 \geqslant \ldots \geqslant a_n$, называют выражение $T_{a_1, a_2, \ldots, a_n}=\sum_{\sigma \in S_n} x_{\sigma(1)}^{a_1} x_{\sigma(2)}^{a_2} \ldots x_{\sigma(n)}^{a_n}$, где суммирование ведётся по множеству $S_n$ всех перестановок множества $\{1,2, \ldots, n\}$.

Пусть даны два упорядоченных набора целых неотрицательных чисел $a_1 \geqslant a_2 \geqslant \ldots \geqslant a_n$ и $b_1 \geqslant b_2 \geqslant$ $\geqslant \ldots \geqslant b_n$. Тогда набор $a_k$ \textit{мажорирует} набор $b_k$, если выполнены следующие условия: $a_1 \geqslant b_1, a_1+$ $+a_2 \geqslant b_1+b_2, \ldots, a_1+\ldots+a_{n-1} \geqslant b_1+\ldots+b_{n-1}, a_1+\ldots+a_n=b_1+\ldots+b_n$. Если набор $a_k$ мажорирует набор $b_k$, то для неотрицательных значений переменных верно \textit{неравенство Мюрхеда}: $T_{a_1, a_2, \ldots, a_n} \geqslant T_{b_1, b_2, \ldots, b_n}$.

\problem{1}
Для положительных чисел $x, y, z$ докажите неравенство

$$
(x+y)(y+z)(z+x)+x y z \leqslant 3\left(x^3+y^3+z^3\right)
$$

\problem{2}
\textbf{Неравенство Шура.} Докажите, что для любых неотрицательных чисел $x, y, z$ и натурального $n \geqslant 3$ выполнено неравенство

$$
T_{n, 0,0}(x, y, z)+T_{n-2,1,1}(x, y, z) \geqslant 2 T_{n-1,1,0}(x, y, z)
$$

\problem{3}
Сумма неотрицательных чисел $x, y$ и $z$ равна 1. Докажите неравенство

$$
0 \leqslant x y+y z+z x-2 x y z \leqslant \frac{7}{27}
$$


\problem{4}
Сумма неотрицательных чисел $x, y$ и $z$ равна 1. Докажите неравенство

$$
1+9 x y z \geqslant 4(x y+y z+z x)
$$

\problem{5}
Для положительных чисел $a, b, c$ докажите неравенство

$$
\frac{(a+b-c)^2}{a b}+\frac{(b+c-a)^2}{b c}+\frac{(c+a-b)^2}{c a} \geqslant 3
$$

\problem{6}
Докажите для положительных $a, b, c$ неравенство

$$
\frac{c}{(a+b)^2}+\frac{a}{(b+c)^2}+\frac{b}{(a+c)^2} \geqslant \frac{9}{4(a+b+c)}
$$

\problem{7}
Для положительных чисел $a, b, c$ и $d$ докажите неравенство

$$
\sqrt{\frac{a^2+b^2+c^2+d^2}{4}} \geqslant \sqrt[3]{\frac{a b c+b c d+c d a+d a b}{4}}
$$

\problem{8}
Положительные числа $a, b, c$ таковы, что $a b+b c+c a=a+b+c$. Докажите, что

$$
a+b+c+1 \geqslant 4 a b c
$$